% Zumdahl Chapter 14 free-response set -- 2 long (10) + 3 short (4) = 32 pts.
\documentclass{apchem-exam}

\apcunitword{Chapter}
\apcunit{14}
\apcunittitle{Acids and Bases}
\apcdoctitle{Free-Response Set}
\apcsubtitle{Zumdahl \S14.1--\S14.9 \textbullet\ 5 questions \textbullet\ 32 points \textbullet\ 50 minutes}

\begin{document}
\apctitleblock

\begin{apcnote}[Scope --- one CED exclusion applies]
Covers \textbf{CED 8.1} (\S14.1--14.2), \textbf{8.2} (\S14.3--14.4,
\S14.6), \textbf{8.3} (\S14.5--14.7), \textbf{8.6} (\S14.9) and
\textbf{8.7} (\S14.8).

The CED exclusion on topic 4.8 removes \textbf{Lewis acid--base
concepts}. Zumdahl develops the Lewis model in \S14.11; it is not
assessed and no question below uses it. Br\o{}nsted--Lowry --- proton
donor and acceptor --- is the model the exam expects throughout.

$K_w = 1.0\times10^{-14}$ at \SI{25}{\celsius}.
\end{apcnote}

\examsection{II}{Free Response}{5 questions \textbullet\ 32 points \textbullet\ 50 minutes}{\calculatorok}

% =====================================================================
\frq{long}{10}{Zumdahl \S14.5 \textbullet\ CED 8.3 \textbullet\ the pH of a weak acid}

Acetic acid, \ce{CH3COOH}, has $K_a = 1.8\times10^{-5}$. Consider a
\SI{0.100}{\Molar} solution.

\begin{parts}
  \item Write the equation for its ionization in water and the $K_a$
        expression. \answerlines[3]{\ce{CH3COOH(aq) + H2O(l) <=>
        CH3COO^-(aq) + H3O+(aq)}

        $K_a = \dfrac{[\ce{CH3COO^-}][\ce{H3O+}]}{[\ce{CH3COOH}]}$ ---
        water is the solvent and is omitted.}
  \item Calculate the pH of the solution. \workspace[3.6cm]
        \keyonly{\begin{apcanswer}Let $x = [\ce{H3O+}]$ at equilibrium.

        $K_a = \dfrac{x^2}{0.100 - x} = 1.8\times10^{-5}$

        Assume $x \ll 0.100$, so $0.100 - x \approx 0.100$:

        $x = \sqrt{(1.8\times10^{-5})(0.100)} = \num{1.34e-3}$~M

        $\text{pH} = -\log(1.34\times10^{-3}) = \mathbf{2.87}$

        (Solving the quadratic without the approximation gives
        $x = 1.33\times10^{-3}$ and pH \num{2.88}. Either answer is
        correct to the precision the data support.)\end{apcanswer}}
  \item Justify the approximation you used, or explain why none was
        needed. \answerlines[3]{The ionized fraction is
        $1.34\times10^{-3}/0.100 = \SI{1.3}{\percent}$, comfortably below
        the usual \SI{5}{\percent} threshold. So little acid ionizes that
        the equilibrium concentration is indistinguishable from the
        initial one at this precision, and the approximation is
        justified.}
  \item Calculate the percent ionization, and state what happens to it if
        the solution is diluted. \answerlines[3]{Percent ionization
        $= \SI{1.3}{\percent}$.

        On dilution it \textbf{increases}. $K_a$ is fixed, so as the
        concentration falls a \emph{larger fraction} of the acid must
        ionize to satisfy it --- even though the actual
        $[\ce{H3O+}]$ falls and the pH rises.}
\end{parts}

\begin{rubric}
  \rpoint{2}{(a) 1 pt the equilibrium with the double arrow; 1 pt the
    $K_a$ expression with water omitted}
  \rpoint{3}{(b) 1 pt correct $K_a$ setup; 1 pt
    $x = \SI{1.3e-3}{\Molar}$; 1 pt pH \num{2.87} (accept
    \num{2.87}--\num{2.88})}
  \rpoint{2}{(c) 1 pt computing the ionized fraction; 1 pt comparing it
    with the \SI{5}{\percent} rule}
  \rpoint{3}{(d) 1 pt \SI{1.3}{\percent}; 1 pt increases on dilution;
    1 pt the reasoning that $K_a$ is constant. Saying pH falls on dilution
    earns 0 for the last point}
\end{rubric}

% =====================================================================
\frq{long}{10}{Zumdahl \S14.2--\S14.4, \S14.9 \textbullet\ CED 8.1, 8.2, 8.6 \textbullet\ strength, pH and structure}

\begin{parts}
  \item Distinguish between a \emph{strong} acid and a
        \emph{concentrated} acid. \answerlines[3]{\textbf{Strength} is
        about the extent of ionization: a strong acid ionizes essentially
        completely in water. \textbf{Concentration} is about how much acid
        was dissolved per litre. A dilute solution of a strong acid and a
        concentrated solution of a weak acid are both entirely possible.}
  \item Calculate the pH of \SI{0.0100}{\Molar} \ce{HCl} and of
        \SI{0.0100}{\Molar} \ce{NaOH}. \workspace[3cm]
        \keyonly{\begin{apcanswer}\ce{HCl} is a strong acid, fully
        ionized, so $[\ce{H3O+}] = \SI{0.0100}{\Molar}$ and
        $\text{pH} = \mathbf{2.00}$.

        \ce{NaOH} is a strong base, so
        $[\ce{OH-}] = \SI{0.0100}{\Molar}$ and $\text{pOH} = 2.00$.

        $\text{pH} = 14.00 - 2.00 = \mathbf{12.00}$\end{apcanswer}}
  \item Explain why \ce{HF} is a weak acid while \ce{HCl} is strong, even
        though fluorine is the more electronegative element.
        \workspace[3cm]
        \keyonly{\begin{apcanswer}Acid strength here is governed by
        \textbf{bond strength}, not by electronegativity. Fluorine is
        small, so the \ce{H-F} bond is short and very strong --- much
        harder to break than the longer, weaker \ce{H-Cl} bond.

        Ionization requires that bond to break. \ce{HCl} gives up its
        proton readily and ionizes completely; \ce{HF} holds onto its
        proton, so only a small fraction ionizes. Down the halogen group
        the bond lengthens and weakens, and acid strength
        \emph{increases}.\end{apcanswer}}
  \item Oxoacid strength increases along \ce{HClO}, \ce{HClO2},
        \ce{HClO3}, \ce{HClO4}. Explain the trend.
        \answerlines[3]{Each additional oxygen is highly electronegative
        and withdraws electron density from the \ce{O-H} bond, weakening
        it and making the proton easier to release. The extra oxygens also
        spread the negative charge of the resulting anion over more atoms,
        so the anion is more able to accommodate it and the ionization
        proceeds further.}
\end{parts}

\begin{rubric}
  \rpoint{2}{(a) 1 pt strength as extent of ionization; 1 pt concentration
    as amount per volume}
  \rpoint{3}{(b) 1 pt pH \num{2.00}; 1 pt pOH \num{2.00}; 1 pt pH
    \num{12.00}. Answering pH \num{2.00} for the base earns 0 for that
    half}
  \rpoint{3}{(c) 1 pt identifying bond strength as the deciding factor;
    1 pt \ce{H-F} short and strong; 1 pt so \ce{HF} ionizes only slightly.
    An answer that argues from electronegativity alone predicts the wrong
    order and earns \textbf{0}}
  \rpoint{2}{(d) 1 pt electron withdrawal weakening the \ce{O-H} bond;
    1 pt charge spread over more oxygens in the anion}
\end{rubric}

% =====================================================================
\frq{short}{4}{Zumdahl \S14.3 \textbullet\ CED 8.2 \textbullet\ pH, pOH and \emph{K}\textsubscript{w}}

A solution at \SI{25}{\celsius} has
$[\ce{OH-}] = \SI{2.5e-4}{\Molar}$.

\begin{parts}
  \item Calculate the pOH, the $[\ce{H3O+}]$ and the pH.
        \workspace[3cm]
        \keyonly{\begin{apcanswer}$\text{pOH} = -\log(2.5\times10^{-4})
        = \mathbf{3.60}$

        $\text{pH} = 14.00 - 3.60 = \mathbf{10.40}$

        $[\ce{H3O+}] = \dfrac{1.0\times10^{-14}}{2.5\times10^{-4}}
        = \mathbf{4.0\times10^{-11}}$~M\end{apcanswer}}
  \item State whether the solution is acidic or basic, and how you know.
        \answerlines[2]{\textbf{Basic}: the pH exceeds 7 at
        \SI{25}{\celsius}, and equivalently
        $[\ce{OH-}] > [\ce{H3O+}]$.}
\end{parts}

\begin{rubric}
  \rpoint{3}{(a) 1 pt pOH \num{3.60}; 1 pt pH \num{10.40}; 1 pt
    $[\ce{H3O+}] = \SI{4.0e-11}{\Molar}$}
  \rpoint{1}{(b) basic, with either justification}
\end{rubric}

% =====================================================================
\frq{short}{4}{Zumdahl \S14.1 \textbullet\ CED 8.1 \textbullet\ conjugate acid--base pairs}

\begin{parts}
  \item Identify the two conjugate acid--base pairs in
        \[ \ce{HCO3^-(aq) + H2O(l) <=> H2CO3(aq) + OH^-(aq)} \]
        \answerlines[3]{Pair 1: \ce{H2CO3} (acid) and \ce{HCO3-} (base).
        \quad Pair 2: \ce{H2O} (acid) and \ce{OH-} (base). Each pair
        differs by exactly one proton.}
  \item \ce{HCO3-} can act as either an acid or a base. State the term for
        this and give the equation for it acting as an acid.
        \answerlines[3]{It is \textbf{amphiprotic} (amphoteric).

        \ce{HCO3^-(aq) + H2O(l) <=> CO3^2-(aq) + H3O+(aq)}}
\end{parts}

\begin{rubric}
  \rpoint{2}{(a) 1 pt each correctly matched pair. Pairing \ce{HCO3-} with
    \ce{OH-} earns 0 --- they differ by more than one proton}
  \rpoint{2}{(b) 1 pt amphiprotic/amphoteric; 1 pt a correct equation
    showing proton loss}
\end{rubric}

% =====================================================================
\frq{short}{4}{Zumdahl \S14.8 \textbullet\ CED 8.7 \textbullet\ acid--base properties of salts}

Predict whether an aqueous solution of each salt is acidic, basic or
neutral, and justify.

\begin{parts}
  \item \ce{NaCl} \answerlines[2]{\textbf{Neutral}. \ce{Na+} comes from a
        strong base and \ce{Cl-} from a strong acid, so neither ion reacts
        appreciably with water.}
  \item \ce{NH4Cl} \answerlines[2]{\textbf{Acidic}. \ce{NH4+} is the
        conjugate acid of the weak base ammonia and donates a proton to
        water; \ce{Cl-} is inert.}
  \item \ce{NaCH3COO} \answerlines[2]{\textbf{Basic}. Acetate is the
        conjugate base of a weak acid and accepts a proton from water,
        producing \ce{OH-}; \ce{Na+} is inert.}
\end{parts}

\begin{rubric}
  \rpoint{3}{1 pt each correct prediction}
  \rpoint{1}{a justification naming the parent acid or base in at least
    two cases. ``It contains sodium so it is basic'' earns 0}
\end{rubric}

\examtotal

\end{document}
